% Hay un \newpage


\documentclass[a4paper,11pt]{article}
\usepackage[spanish,activeacute]{babel}
\usepackage[latin1]{inputenc}

\usepackage{multicol}
\usepackage{amssymb, amsmath}
\usepackage{latexsym}
\usepackage{multicol}


\renewcommand{\Re}{\operatorname{Re}}
\renewcommand{\Im}{\operatorname{Im}}

\setlength{\unitlength}{1mm}
\addtolength{\voffset}{-15mm} \addtolength{\textheight}{30mm}
\addtolength{\hoffset}{-15mm} \addtolength{\textwidth}{30mm}


\begin{document}

\begin{flushright}
1${}^o$ de grado en F�sica, curso 2013/14
\end{flushright}



\begin{center}{\bf {\Large \'ALGEBRA I. HOJA 1}}\end{center}

\vskip.5cm

\begin{enumerate}

\item Escribir en coordenadas polares $4 + i, - 3/2 - i/2, -1 + 2
i$.

\item Comprobar que $(1 + i)^{12} = -64$, y $((1-i)/\sqrt2)^{-6} =
- i$.


\item Usando \enspace $e^{ix}=\cos x + i \operatorname{sen} x$,
\enspace donde $x$
es un n�mero real, demostrar que
\newline
$\cos x = (e^{ix} + e^{-ix})/2$ \enspace y \enspace $\operatorname{sen} x = (e^{ix}
- e^{-ix})/(2i)$.

\item Describir el conjunto del plano complejo determinado por las
siguientes  relaciones:


(i) $|z - 2| > |z+ 2|$; \qquad \qquad (ii)   $\Re z + \Im z <1$.

% \medskip (iii)  $|2z| > |1+z^2|.$


\item Demostrar las f�rmulas:

(i) $|z + 1|^2 - |z - 1|^2 = 4\Re z$;
\qquad \qquad
(ii)   $|z + 1|^2 + |z - 1|^2 = 2+ 2\, |z|^2$.


\item  Resolver $z^3 = i$ en coordenadas polares y rectangulares.
(Recordatorio: las coordenadas rectangulares, o cartesianas,
 son de la forma
$a + i b$; las polares, $ r e^{i\alpha}$).

Soluci�n (parcial): $\sqrt{3}/2 + i/2, - \sqrt{3}/2 + i/2, - i$.

\item Resolver $iz + (2 -10 i) z = 3 z + 2i$.

Respuesta: $-9/41  - i/41$.

\item Resolver $z^2 - (3 +i) z + 4 + 3i = 0$.

Respuesta: $2 - i, 1 + 2i$.


\item Resolver:
\textbf{(a)} $z^2= - 8 - 6i$; \quad
\textbf{(b)} $z^2 - 4z + 4 + 2i = 0$; \quad
\textbf{(c)} $z^2 + 2 z + 1 -8i = 0$.  \quad


\item Resolver las ecuaciones

\textbf{(a)}
$
\det
\begin{pmatrix}
1-\lambda & 1 \\
- 2 & 1-\lambda
\end{pmatrix} = 0;
$;
\qquad
\textbf{(b)}
$
\det
\begin{pmatrix}
1-\lambda & 1 \\
- 2 & 3-\lambda
\end{pmatrix} = 0
$.




\item Encontrar todos los autovalores complejos de las
siguientes matrices:

\textbf{(a)}
$
\begin{pmatrix}
3 & 1 \\
- 4 & 3
\end{pmatrix}
$;
\qquad
\textbf{(b)}
$
\begin{pmatrix}
1 & i \\
2 & 1
\end{pmatrix}
$;
\qquad
\textbf{(c)}
$
\begin{pmatrix}
0 & 1 & 0 & 0\\
0 & 0 & 1 & 0\\
0 & 0 & 0 & 1\\
1 & 0 & 0 & 0
\end{pmatrix}
$\; .

\qquad


\item Escribir  $\sum_{n=0}^{99} (i + 1)^n$ en coordenadas
rectangulares y polares. \enspace
Respuesta: $(1 + 2^{50}) i$.


\item Demostrar que para todo polinomio $p(z)$ con coeficientes reales,
 $\overline{p(z)} = p(\overline{ z})$.


\item Hallar un n�mero complejo en forma bin�mica $a + bi$ tal que
$(a + bi)^2 = 1 + i$.

\item Calcular en forma bin�mica y en forma trigonom�trica los siguientes n�meros complejos:
$$\sqrt{1 + i}, \qquad  \sqrt{-2 + 2i}.$$
Comparando las expresiones determinar el valor de $\cos(\pi/8)$ y $\cos(3\pi/8)$.

Comprobar que $\cos^2 (\pi/8)+
\cos^2 (3\pi/8) = 1$. �Por qu�?


%%%%%%%%%%%%%%%%%%%%%%%%%%%%%%%%%%%%%%%%%%%%%%%%%%%%%%%%%%%%%%%%%%%%%%%%%

 \newpage

%%%%%%%%%%%%%%%%%%%%%%%%%%%%%%%%%%%%%%%%%%%%%%%%%%%%%%%%%%%%%%%%%%%%%%%%%



\item Demostrar
\begin{multicols}{2}
\begin{itemize}
\item[a)] $\overline{z_1 + z_2} =\bar z_1 + \bar z_2 $;
\item[b)] $\overline{z_1\cdot z_2} = \bar z_1 \cdot \bar z_2$.
\end{itemize}
\end{multicols}



\item Hallar las ra�ces cuartas de $-i$ y representarlas gr�ficamente.


\item Hallar las siguientes raices:
$$\sqrt[3]{ \frac{\sqrt{3}}{2} + i\frac{1}{2}},\qquad
\sqrt[6]{\frac{1}{2} + i\frac{
\sqrt{3}}{2}}. $$
�C�mo est�n relacionadas entre s�?



\item Comprobar que $(x^{n-1} + x^{n-2} + x^{n-3} +
\cdots + x + 1)(x- 1) = x^n-1$. Deducir las siguientes propiedades.
\begin{itemize}
\item[a)] Las soluciones complejas y las soluciones reales de $x^n + x^{n-1} + x^{n-2} +
\cdots + x + 1 = 0$ son ra�ces $(n + 1)$-�simas de la unidad.
\item[b)] Las ra�ces $(n + 1)$-�simas de la unidad que no coinciden con $1$, son soluciones de la ecuaci�n
$x^n + x^{n-1} + x^{n-2} + \cdots + x + 1 = 0$.
\item[c)] La ecuaci�n $x^4 + x^3 + x^2 + x + 1 = 0$ no tiene ninguna soluci�n real.
\item[d)] Demostrar que la ecuaci�n $x^n + x^{n-1} + x^{n-2} +
\cdots+ x + 1 = 0$ no tiene ninguna soluci�n real si $n$ es par y
tiene exactamente una soluci�n real si $n$ es impar. �Cu�l es la soluci�n real si $n$ es impar?
\end{itemize}



\item Deducir de la f�rmula de De Moivre
$(\cos \alpha + i \sen\alpha)^n = \cos (n\alpha) + i \sen(n\alpha)$:
\begin{itemize}
\item[a)] las f�rmulas del coseno del �ngulo triple y del seno del �ngulo triple;
\item[b)] las f�rmulas an�logas para el �ngulo qu�ntuple.
\end{itemize}


\item Hallar $\cos(\pi/12)$ calculando la ra�z de $e^{i\pi/6}$.

\end{enumerate}

\end{document}
